%% 第5巻 巻末（こたえ・記号カード・定数と単位・さくいん・奥付）
\makeatletter
%% ---------------- こたえ ----------------
\bkAnswerPages

%% ---------------- 記号カード ----------------
\bkKigoPages{%
\bkkigorow{$\Delta$}{デルタ}{変化量。「あと$-$まえ」}{2-3}%
\bkkigorow{$f(x)$}{エフ・エックス}{関数 $f$ に $x$ を入れた値}{3-3}%
\bkkigorow{$\theta$}{シータ}{角の大きさ（ラジアン）}{4-2}%
\bkkigorow{$\omega$}{オメガ}{角速度。1秒に回る角（rad/s）}{4-8}%
}

%% 記号カードの2ページ目：ギリシャ文字（第6回：一覧のつづきのすぐあとに置く）
\def\bkbkhead#1{\par\addvspace{4mm}\noindent{\sffamily\bfseries\large\color{bkTheme}#1}\par\nobreak\vspace{2mm}}
\bkbkhead{ギリシャ文字の読み方}
{\small 工学の本では、ギリシャ文字がたくさん使われます。大文字・小文字・読み方と、よく使われる意味です。}\par\vspace{2mm}
\begin{bkcenter}
\footnotesize
\begin{tblr}{colspec={Q[c,m,wd=7mm] Q[c,m,wd=7mm] Q[l,m,wd=17mm] Q[l,m,wd=30mm] Q[c,m,wd=7mm] Q[c,m,wd=7mm] Q[l,m,wd=17mm] Q[l,m,wd=30mm]}, row{1}={bg=bkVol!14, font=\sffamily\bfseries\footnotesize}, hlines={0.4pt,bkLine}, vline{5}={0.8pt,bkVol}, rowsep=0.8pt, colsep=2.5pt}
大 & 小 & 読み方 & よく使う意味 & 大 & 小 & 読み方 & よく使う意味 \\
$\mathrm{A}$ & $\alpha$ & アルファ & 角、角加速度 & $\mathrm{N}$ & $\nu$ & ニュー & 振動数 \\
$\mathrm{B}$ & $\beta$ & ベータ & 角 & $\Xi$ & $\xi$ & クシー & （変数） \\
$\Gamma$ & $\gamma$ & ガンマ & （係数） & $\mathrm{O}$ & $o$ & オミクロン & ― \\
$\Delta$ & $\delta$ & デルタ & 変化量、小さな量 & $\Pi$ & $\pi$ & パイ & 円周率 \\
$\mathrm{E}$ & $\varepsilon$ & イプシロン & 小さな数、誘電率 & $\mathrm{P}$ & $\rho$ & ロー & 密度 \\
$\mathrm{Z}$ & $\zeta$ & ゼータ & 減衰の割合 & $\Sigma$ & $\sigma$ & シグマ & 和、ばらつき \\
$\mathrm{H}$ & $\eta$ & イータ & 効率 & $\mathrm{T}$ & $\tau$ & タウ & 時定数、トルク \\
$\Theta$ & $\theta$ & シータ & 角 & $\Upsilon$ & $\upsilon$ & ウプシロン & ― \\
$\mathrm{I}$ & $\iota$ & イオタ & ― & $\Phi$ & $\phi$ & ファイ & 位相、磁束 \\
$\mathrm{K}$ & $\kappa$ & カッパ & （係数） & $\mathrm{X}$ & $\chi$ & カイ & ― \\
$\Lambda$ & $\lambda$ & ラムダ & 波長 & $\Psi$ & $\psi$ & プサイ & （角） \\
$\mathrm{M}$ & $\mu$ & ミュー & 摩擦係数、マイクロ & $\Omega$ & $\omega$ & オメガ & 角速度（$\Omega$ は抵抗の単位） \\
\end{tblr}
\end{bkcenter}
{\footnotesize 大文字が英語のアルファベットと同じ形の文字（A、B、E など）は、記号としてはほとんど使いません。$\varepsilon$ は $\epsilon$ と書くこともあります。}

\bkbkhead{大きな数・小さな数を表す記号（接頭語）}
\begin{bkcenter}
\small
\begin{tblr}{colspec={Q[l,m,wd=14mm] *{6}{Q[c,m,wd=18mm]}}, column{1}={bg=bkVol!14, font=\sffamily\bfseries\footnotesize}, hlines={0.4pt,bkLine}, rowsep=1.2pt, colsep=2pt}
記号 & G & M & k & m & $\upmu$ & n \\
読み方 & ギガ & メガ & キロ & ミリ & マイクロ & ナノ \\
倍 & $10^9$ & $10^6$ & $10^3$ & $10^{-3}$ & $10^{-6}$ & $10^{-9}$ \\
例 & GHz & MW & km、kWh & mm、mA & $\upmu$m、$\upmu$F & nm \\
\end{tblr}
\end{bkcenter}

%% 第6回：記号カードの3ページ目。ノートに書くとき、ほかの字とまぎれやすい記号
\clearpage
\bkbkhead{記号を書くときのちゅうい}
{\small この本では、問題をノートに解きます。手で書くと、ほかの字とまぎれやすい記号があります。書くときのコツです。}\par\vspace{2mm}
\begin{bkcenter}
\begin{tblr}{colspec={Q[c,m,wd=22mm] Q[l,m,wd=34mm] Q[l,m,wd=68mm]}, row{1}={bg=bkVol!14, font=\sffamily\bfseries\footnotesize}, hlines={0.4pt,bkLine}, rowsep=3pt, cells={font=\small}}
記号 & まぎれやすいもの & 書くときのコツ \\
$\displaystyle\int_a^b$ & S、$f$ & 上下に長い、細い S の形。$a$（下端）は右下、$b$（上端）は右上に小さく書く \\
$\displaystyle\sum_{k=1}^{n}$ & E & 上と下の横線を同じ長さにし、まん中をとがらせる。$k=1$ は下、$n$ は上 \\
$\Delta x$ & 三角形の記号 △、$\Delta\times x$ & $\Delta x$ で1つの記号（「$x$ の変化量」）。$\Delta$ と $x$ をかけるのではない \\
$\dd x$、$\dfrac{\dd y}{\dd x}$ & 文字の $d$ をかけたもの & $\dd x$ もひとまとまり。$\dfrac{\dd y}{\dd x}$ は分数の形をした1つの記号（$d$ で約分しない） \\
$\theta$ & 0、O & たて長の丸の、まん中に横線を1本 \\
$\omega$ & $w$ & 下を丸く2つ。$w$ のようにとがらせない \\
$f'(x)$、$f''(x)$ & 1、$f^1$ & 右かたに小さく短い線。2回微分は2本 \\
$e$ & 文字の $e$ & 2.718… の決まった数。$e^x$ の $e$ は数なので、ほかの文字に変えない \\
$\infty$ & 8、$\propto$ & 8 を横にたおした、閉じた形 \\
$\displaystyle\lim_{n\to\infty}$ & ― & $n\to\infty$ は lim の真下に小さく書く \\
\end{tblr}
\end{bkcenter}


%% ---------------- 定数・単位・公式 ----------------
\bkbacktitle{よく使う定数・単位・公式}
\def\bkbkhead#1{\par\addvspace{4mm}\noindent{\sffamily\bfseries\large\color{bkTheme}#1}\par\nobreak\vspace{2mm}}
\bksubsection{1}{この巻で使う定数}
\begin{bkcenter}
\small
\begin{tblr}{colspec={Q[c,m,wd=14mm] Q[l,m,wd=34mm] Q[l,m,wd=34mm] Q[c,m,wd=14mm]}, row{1}={bg=bkVol!14, font=\sffamily\bfseries\footnotesize}, hlines={0.4pt,bkLine}, rowsep=1.6pt}
記号 & 名前 & 値 & 章 \\
$g$ & 重力加速度 & \qty{9.8}{m/s^2} & 5-5 \\
$\pi$ & 円周率 & 3.14159… & 4-2 \\
$e$ & ネイピア数 & 2.71828… & 5-14 \\
$\ln 2$ & 2の自然対数 & 0.693… & 5-14 \\
$\ln 10$ & 10の自然対数 & 2.303… & 5-14 \\
$\sqrt2$ & 2の平方根 & 1.414… & 2-8 \\
\end{tblr}
\end{bkcenter}

\bksubsection{2}{この巻で使う単位}
\begin{bkcenter}
\small
\begin{tblr}{colspec={Q[l,m,wd=27mm] Q[l,m,wd=27mm] Q[l,m,wd=56.5mm] Q[c,m,wd=12mm]}, row{1}={bg=bkVol!14, font=\sffamily\bfseries\footnotesize}, hlines={0.4pt,bkLine}, rowsep=1.6pt}
量 & 単位 & 意味・ほかの単位との関係 & 章 \\
速度 & m/s & 1秒に進む距離。$\qty{1}{m/s}=\qty{3.6}{km/h}$ & 5-1 \\
加速度 & \unit{m/s^2} & 1秒に変わる速度 & 5-5 \\
力 & N（ニュートン） & $\qty{1}{N}=\qty{1}{kg.m/s^2}$ & 5-12 \\
仕事・エネルギー & J（ジュール） & $\qty{1}{J}=\qty{1}{N.m}$ & 5-12 \\
仕事率・電力 & W（ワット） & $\qty{1}{W}=\qty{1}{J/s}$ & 5-12 \\
電力量 & kWh & $\qty{1}{kWh}=\qty{3.6e6}{J}$ & 5-12 \\
電気の量 & C（クーロン） & \qty{1}{A}の電流が1秒流れたときの量 & 5-12 \\
流量 & L/s & 1秒に流れる水の量 & 5-12 \\
角 & rad（ラジアン） & 弧の長さ÷半径。$180^\circ=\pi$\,rad & 4-2 \\
角速度 & rad/s & 1秒に回る角 & 5-15 \\
\end{tblr}
\end{bkcenter}


\clearpage
\bksubsection{3}{微分と積分の表}
\begin{bkcenter}
\small
\begin{tblr}{colspec={Q[c,m,wd=26mm] Q[c,m,wd=32mm] Q[c,m,wd=44mm] Q[c,m,wd=18mm]}, row{1}={bg=bkVol!14, font=\sffamily\bfseries\footnotesize}, hlines={0.4pt,bkLine}, rowsep=2pt}
関数 & 微分すると & 積分すると（$+C$ を省略） & 章 \\
$c$（定数） & $0$ & $cx$ & 5-4 \\
$x^n$ & $nx^{n-1}$ & $\dfrac{x^{n+1}}{n+1}$（$n\neq-1$） & 5-4・5-11 \\
$e^{ax}$ & $ae^{ax}$ & $\dfrac{1}{a}e^{ax}$ & 5-14 \\
$\dfrac{1}{x}$ & $-\dfrac{1}{x^2}$ & $\ln x$（$x>0$） & 5-14 \\
$\ln x$ & $\dfrac{1}{x}$ & $x\ln x-x$ & 5-14・5-16 \\
$\sin\omega x$ & $\omega\cos\omega x$ & $-\dfrac{1}{\omega}\cos\omega x$ & 5-15 \\
$\cos\omega x$ & $-\omega\sin\omega x$ & $\dfrac{1}{\omega}\sin\omega x$ & 5-15 \\
\end{tblr}
\end{bkcenter}

\bksubsection{4}{計算のきまり}
\begin{bkcenter}
\small
\begin{tblr}{colspec={Q[l,m,wd=28mm] Q[l,m,wd=82.5mm] Q[c,m,wd=16mm]}, row{1}={bg=bkVol!14, font=\sffamily\bfseries\footnotesize}, hlines={0.4pt,bkLine}, rowsep=2pt, column{1}={font=\sffamily}}
きまり & 式 & 章 \\
積の微分 & $(fg)'=f'g+fg'$ & 5-6 \\
合成関数の微分 & $\diff{y}{x}=\diff{y}{u}\cdot\diff{u}{x}$ & 5-6 \\
1次近似 & $f(a+h)\approx f(a)+f'(a)\,h$ & 5-8 \\
微積分の基本定理 & $F'=f$ なら $\displaystyle\int_a^bf\,\dd x=F(b)-F(a)$ & 5-11 \\
置換積分 & $u=g(x)$ とおき、$\dd u=g'(x)\,\dd x$ で書きかえる。定積分は、はんいも $u$ に & 5-16 \\
部分積分 & $\displaystyle\int fg'\,\dd x=fg-\int f'g\,\dd x$ & 5-16 \\
2倍角から & $\sin^2\theta=\dfrac{1-\cos2\theta}{2}$、\quad $\cos^2\theta=\dfrac{1+\cos2\theta}{2}$ & 4-9・5-15 \\
テイラー展開 & $f(x)=f(0)+f'(0)\,x+\dfrac{f''(0)}{2!}x^2+\cdots$ & 5-17 \\
 & $e^x=1+x+\dfrac{x^2}{2!}+\cdots$、\ $\sin x=x-\dfrac{x^3}{3!}+\cdots$、\ $\cos x=1-\dfrac{x^2}{2!}+\cdots$ & 5-17 \\
\end{tblr}
\end{bkcenter}

\clearpage
\bksubsection{5}{よく出る関数のグラフ}
{\small この巻で微分・積分した関数のグラフの形です。グラフの形を思いうかべると、微分の符号（ふえる・へる）や、積分の値のおよそがわかります。}\par\vspace{2mm}
\begin{bkcenter}
\def\bkgr#1#2#3{\begin{minipage}[t]{0.32\linewidth}\centering
\begin{tikzpicture}
\begin{axis}[bkplot, width=46mm, height=38mm, #1]
#2
\end{axis}
\end{tikzpicture}\par\vspace{1mm}
{\footnotesize #3}
\end{minipage}}
\bkgr{xmin=-2.2, xmax=2.2, ymin=-1, ymax=4.5, xtick={-2,0,2}, ytick={2,4}}{\addplot[bkcurve, domain=-2.1:2.1, samples=60] {x^2};}{$y=x^2$\quad 微分 $2x$\\$x=0$ で傾き0（谷）}\hfill
\bkgr{xmin=-1.8, xmax=1.8, ymin=-4, ymax=4, xtick={-1,1}, ytick={-2,2}}{\addplot[bkcurve, domain=-1.58:1.58, samples=60] {x^3};}{$y=x^3$\quad 微分 $3x^2$\\いつも傾きは0以上}\hfill
\bkgr{xmin=0, xmax=4.2, ymin=0, ymax=4.5, xtick={1,2,3,4}, ytick={1,2,3,4}}{\addplot[bkcurve, domain=0.23:4.1, samples=60] {1/x};}{$y=\dfrac{1}{x}$（$x>0$）\quad 微分 $-\dfrac{1}{x^2}$\\積分すると $\ln x$}

\vspace{4mm}
\bkgr{xmin=-2.2, xmax=2.2, ymin=0, ymax=7.5, xtick={-2,0,2}, ytick={1,4}}{\addplot[bkcurve, domain=-2.1:2, samples=60] {exp(x)}; \addplot[bkcurve2, domain=-2:2.1, samples=60] {exp(-x)};}{$y=e^x$（青）・$e^{-x}$\mbox{（オレンジ）}\\微分しても形が変わらない}\hfill
\bkgr{xmin=0, xmax=6.2, ymin=-2.2, ymax=2.2, xtick={1,2.718}, xticklabels={1,$e$}, ytick={-1,1}}{\addplot[bkcurve, domain=0.12:6, samples=80] {ln(x)};}{$y=\ln x$\quad 微分 $\dfrac{1}{x}$\\$(1,\,0)$ を通り、ゆっくりふえる}\hfill
\bkgr{xmin=0, xmax=6.6, ymin=-1.3, ymax=1.3, xtick={1.5708,3.1416,6.2832}, xticklabels={$\frac{\pi}{2}$,$\pi$,$2\pi$}, ytick={-1,1}}{\addplot[bkcurve, domain=0:6.4, samples=80] {sin(deg(x))}; \addplot[bkcurve2, domain=0:6.4, samples=80] {cos(deg(x))};}{$y=\sin x$（青）・$\cos x$\mbox{（オレンジ）}\\微分すると $\frac{1}{4}$ 周期進む}
\end{bkcenter}

\vspace{4mm}
\begin{bkcenter}
\small
\begin{tblr}{colspec={Q[l,m,wd=40mm] Q[l,m,wd=88mm]}, row{1}={bg=bkVol!14, font=\sffamily\bfseries\footnotesize}, hlines={0.4pt,bkLine}, rowsep=1.6pt}
グラフのようす & 微分・積分では \\
右上がり（ふえる） & 微分した値（傾き）が正 \\
右下がり（へる） & 微分した値が負 \\
山・谷（いちばん高い・低い所） & 微分した値が0（前後で符号が変わる） \\
下に凸・上に凸 & 2回微分した値が正・負 \\
横軸より上・下 & 定積分（面積）が正・負 \\
\end{tblr}
\end{bkcenter}

\clearpage
\bksubsection{6}{表計算の関数}
{\small この巻の「やってみよう」で使った、表計算ソフトの関数です。角はラジアンで入れます。}\par\vspace{2mm}
\begin{bkcenter}
\small
\begin{tblr}{colspec={Q[l,m,wd=40mm] Q[l,m,wd=53.5mm] Q[l,m,wd=32mm]}, row{1}={bg=bkVol!14, font=\sffamily\bfseries\footnotesize}, hlines={0.4pt,bkLine}, rowsep=2pt}
書き方 & 意味 & 例 \\
\texttt{=A2\textasciicircum 3} & A2の3乗 & 2なら8 \\
\texttt{=SQRT(A2)} & A2の平方根 & 9なら3 \\
\texttt{=EXP(A2)} & $e$ のA2乗 & 1なら2.71828… \\
\texttt{=LN(A2)} & A2の自然対数 $\ln$ & 1なら0 \\
\texttt{=SIN(A2)}\quad\texttt{=COS(A2)} & $\sin$、$\cos$（A2はラジアン） & 0なら0と1 \\
\texttt{=PI()} & 円周率 $\pi$ & 3.14159… \\
\texttt{=RADIANS(A2)} & 度をラジアンに直す & 180なら3.14159… \\
\texttt{=SUM(B2:B7)} & B2からB7までの合計（$\Sigma$） & 面積のたし算（5-10） \\
\texttt{=B2+(C2*0.1)} & となりのセルを使って計算 & オイラー法（5-8・5-13） \\
\end{tblr}
\end{bkcenter}

\bksubsection{7}{表計算で計算するときの注意}
\begin{itemize}[leftmargin=1.4\zw, itemsep=2pt, topsep=2pt, label={\color{bkTheme}●}]
\item 角は、度ではなくラジアンで入れます。度の値は \texttt{RADIANS} でなおします（5-15）。
\item $h$ を小さくしすぎると、計算機のけた数の限界で、かえって答えがずれることがあります。$h=0.001$ くらいが目安です（5-3）。
\item 式をとなりのセルにコピーすると、セルの番地がいっしょにずれます。ずらしたくない番地は \texttt{\$B\$1} のように \texttt{\$} をつけます。
\item 表計算の答えは、いつも「近似値」です。式で求めた答えとくらべて、どれくらい合っているかを確かめるのが、この本のやり方です。
\end{itemize}

%% ---------------- チェック表 ----------------
\bkbacktitle{できるようになったことチェック表}
{\small 各章の「ふりかえり」を集めた表です。章を学び終えたら1回め、次の巻に進む前に2回めのチェックをしよう。できなかった所は、ページにもどって読みなおします。}\par\vspace{2mm}
%% 自動生成（tools/mkchecklist.py）。各章の「ふりかえり」の項目を集めたもの
\begingroup\footnotesize
\def\bkcb{\tikz[baseline=-0.5ex]\draw[bkTheme, line width=0.7pt, rounded corners=0.5pt] (0,-0.3em) rectangle (0.8em,0.5em);}
\begin{longtblr}{colspec={Q[c,m,wd=9mm] X[l,m] Q[c,m,wd=10mm] Q[c,m,wd=9mm] Q[c,m,wd=9mm]}, width=\linewidth, rowhead=1, colsep=3pt,
  row{1}={bg=bkVol!14, font=\sffamily\bfseries\scriptsize}, hlines={0.4pt,bkLine}, rowsep=0.7pt}
章 & できるようになったこと & ページ & 1回め & 2回め \\
{\sffamily\bfseries\color{bkTheme}5-1} & 平均の速さを $\tfrac{\Delta x}{\Delta t}$ で求められる。 & \pageref{bk:5-1:heikin} & \bkcb & \bkcb \\
 & 平均の速さが、2点を結ぶ直線の傾きだと説明できる。 & \pageref{bk:5-1:katamuki} & \bkcb & \bkcb \\
 & あいだを短くして、瞬間の速さを見積もれる。 & \pageref{bk:5-1:shunkan} & \bkcb & \bkcb \\
{\sffamily\bfseries\color{bkTheme}5-2} & $\Delta t$ を小さくして、平均の速さを表で計算できる。 & \pageref{bk:5-2:hyou} & \bkcb & \bkcb \\
 & 平均の速さを $h$ の式で表せる。 & \pageref{bk:5-2:h} & \bkcb & \bkcb \\
 & $x=at^2$ の瞬間の速さを求められる。 & \pageref{bk:5-2:a} & \bkcb & \bkcb \\
{\sffamily\bfseries\color{bkTheme}5-3} & $\lim$ の記号を読み、極限の値を求められる。 & \pageref{bk:5-3:lim} & \bkcb & \bkcb \\
 & 微分係数を、接線の傾きとして求められる。 & \pageref{bk:5-3:bibunkeisu} & \bkcb & \bkcb \\
 & 導関数を求め、接線の式を書ける。 & \pageref{bk:5-3:shiki} & \bkcb & \bkcb \\
{\sffamily\bfseries\color{bkTheme}5-4} & 表の結果から、微分のきまりを見つけられる。 & \pageref{bk:5-4:kimari0} & \bkcb & \bkcb \\
 & $(x^n)'=nx^{n-1}$ を使って微分できる。 & \pageref{bk:5-4:kimari} & \bkcb & \bkcb \\
 & 定数倍と和のきまりで、多項式を微分できる。 & \pageref{bk:5-4:wa} & \bkcb & \bkcb \\
{\sffamily\bfseries\color{bkTheme}5-5} & 速度と速さのちがいを説明できる。 & \pageref{bk:5-5:sokudo} & \bkcb & \bkcb \\
 & 位置を2回微分して、加速度を求められる。 & \pageref{bk:5-5:nikai} & \bkcb & \bkcb \\
 & $\dot{x}$、$\ddot{x}$、$\tfrac{\dd^2x}{\dd t^2}$ を読める。 & \pageref{bk:5-5:nikai} & \bkcb & \bkcb \\
{\sffamily\bfseries\color{bkTheme}5-6} & 積の微分を、長方形の面積で説明できる。 & \pageref{bk:5-6:seki} & \bkcb & \bkcb \\
 & 合成関数の内側と外側を見分けられる。 & \pageref{bk:5-6:gousei} & \bkcb & \bkcb \\
 & 合成関数の微分を使って計算できる。 & \pageref{bk:5-6:chain} & \bkcb & \bkcb \\
{\sffamily\bfseries\color{bkTheme}5-7} & 導関数の符号から、上がる・下がるを判断できる。 & \pageref{bk:5-7:fugou} & \bkcb & \bkcb \\
 & 増減表をかいて、極大・極小を求められる。 & \pageref{bk:5-7:zougen} & \bkcb & \bkcb \\
 & 範囲のはしもくらべて、最大・最小を求められる。 & \pageref{bk:5-7:saidai} & \bkcb & \bkcb \\
{\sffamily\bfseries\color{bkTheme}5-8} & 次の位置を、今の位置＋速度×$\Delta t$ で予想できる。 & \pageref{bk:5-8:tsugi} & \bkcb & \bkcb \\
 & 接線を使って、1次近似ができる。 & \pageref{bk:5-8:sessen} & \bkcb & \bkcb \\
 & 測定のずれが、計算にひびく大きさを見積もれる。 & \pageref{bk:5-8:gosa} & \bkcb & \bkcb \\
{\sffamily\bfseries\color{bkTheme}5-9} & 速さが変わるときも、細かく分けて、かけて、たして距離を求められる。 & \pageref{bk:5-9:wakeru} & \bkcb & \bkcb \\
 & 長方形・三角形・台形の面積で、進んだ距離がわかる。 & \pageref{bk:5-9:sankaku} & \bkcb & \bkcb \\
 & マイナスの面積の意味（変位と道のり）を説明できる。 & \pageref{bk:5-9:minus} & \bkcb & \bkcb \\
{\sffamily\bfseries\color{bkTheme}5-10} & 細い長方形の和で、面積を近似できる。 & \pageref{bk:5-10:hasami} & \bkcb & \bkcb \\
 & $\Sigma$ の記号を読み書きできる。 & \pageref{bk:5-10:sigma} & \bkcb & \bkcb \\
 & 定積分の意味を説明し、$x^2$ の定積分を求められる。 & \pageref{bk:5-10:teisekibun} & \bkcb & \bkcb \\
{\sffamily\bfseries\color{bkTheme}5-11} & 面積のふえる速さが高さであることを説明できる。 & \pageref{bk:5-11:kihon} & \bkcb & \bkcb \\
 & 不定積分を $+C$ をつけて書ける。 & \pageref{bk:5-11:genshi} & \bkcb & \bkcb \\
 & 定積分を $F(b)-F(a)$ で計算できる。 & \pageref{bk:5-11:teisekibun} & \bkcb & \bkcb \\
{\sffamily\bfseries\color{bkTheme}5-12} & 変わる力の仕事を、積分で求められる。 & \pageref{bk:5-12:sekibun} & \bkcb & \bkcb \\
 & 流量や電力から、たまった量を求められる。 & \pageref{bk:5-12:ryuuryou} & \bkcb & \bkcb \\
 & 積分を使って、平均の値を求められる。 & \pageref{bk:5-12:heikin} & \bkcb & \bkcb \\
{\sffamily\bfseries\color{bkTheme}5-13} & 初期条件から、積分定数を決められる。 & \pageref{bk:5-13:shoki} & \bkcb & \bkcb \\
 & 加速度から、速度と位置を求められる。 & \pageref{bk:5-13:toukasoku} & \bkcb & \bkcb \\
 & 「いまの量＝はじめの量＋変化の積分」を使える。 & \pageref{bk:5-13:teisekibun} & \bkcb & \bkcb \\
{\sffamily\bfseries\color{bkTheme}5-14} & $e$ が $(1+\frac{1}{n})^n$ の近づく先であることを確かめられる。 & \pageref{bk:5-14:n} & \bkcb & \bkcb \\
 & $(e^{ax})'=ae^{ax}$ を使える。 & \pageref{bk:5-14:eax} & \bkcb & \bkcb \\
 & $\ln x$ を微分・積分できる。 & \pageref{bk:5-14:lnbibun} & \bkcb & \bkcb \\
{\sffamily\bfseries\color{bkTheme}5-15} & 円運動の速度で、$\sin$・$\cos$ の微分を説明できる。 & \pageref{bk:5-15:en} & \bkcb & \bkcb \\
 & $\sin\omega t$、$\cos\omega t$ を微分・積分できる。 & \pageref{bk:5-15:omega} & \bkcb & \bkcb \\
 & 単振動の $\ddot{x}=-\omega^2x$ を確かめられる。 & \pageref{bk:5-15:nikai} & \bkcb & \bkcb \\
 & $\sin^2$ を積分して、交流の実効値を求められる。 & \pageref{bk:5-15:jikko} & \bkcb & \bkcb \\
{\sffamily\bfseries\color{bkTheme}5-16} & 置換積分で、$\dd u$ への書きかえができる。 & \pageref{bk:5-16:chikan} & \bkcb & \bkcb \\
 & 定積分の置換で、はんいを書きかえられる。 & \pageref{bk:5-16:hanni} & \bkcb & \bkcb \\
 & 部分積分で、$f$ と $g'$ を選んで計算できる。 & \pageref{bk:5-16:bubun} & \bkcb & \bkcb \\
{\sffamily\bfseries\color{bkTheme}5-17} & 2次近似で、曲がり方まで合わせる意味を説明できる。 & \pageref{bk:5-17:niji} & \bkcb & \bkcb \\
 & $e^x$・$\sin x$・$\cos x$ の展開を書き、値を計算できる。 & \pageref{bk:5-17:sin} & \bkcb & \bkcb \\
 & $\sin\theta\approx\theta$ から、振り子の周期を求められる。 & \pageref{bk:5-17:furiko} & \bkcb & \bkcb \\
{\sffamily\bfseries\color{bkTheme}5-18} & ふるまい表を使って、微分・積分できる。 & \pageref{bk:5-18:hyo} & \bkcb & \bkcb \\
 & 「自分にもどる」関数を答えられる。 & \pageref{bk:5-18:modoru} & \bkcb & \bkcb \\
 & はじめの状態から、式の中の数を決められる。 & \pageref{bk:5-18:shoki} & \bkcb & \bkcb \\
\end{longtblr}
\endgroup



%% ---------------- つぎの巻へ（第6回で追加：第6巻の予告と、この巻の道具がこの先どこで使われるか） ----------------
\bkbacktitle{つぎの巻へ}
\newcommand\bknextpill[1]{\bk@pill{#1}}% 第11回：章の番号の札は、その章の巻の色（5-○は緑、6-○は青緑）
\def\bknextrow{1}
\newcommand\bknextcard[4]{%
  \begin{tcolorbox}[enhanced, nobeforeafter, box align=top, width=0.485\linewidth, colframe=bkVol6, colback=white, boxrule=0.9pt, arc=2.5mm,
    left=3mm, right=3mm, top=2mm, bottom=2.5mm, equal height group=bknext\bknextrow]
  \bknextpill{#1}\enspace{\sffamily\bfseries\small #2}\par\vspace{1.5mm}
  {\sffamily\bfseries\small\color{bkVol6!80!black}#3}\par\vspace{2mm}
  {\footnotesize\setlength{\parindent}{0pt}{\sffamily\bfseries\color{bkInk}使う道具}\enspace #4}
  \end{tcolorbox}}
微分と積分の道具がそろいました。第6巻「未来を計算する」では、この2つを組み合わせて、「次の瞬間を決めるルール」（微分方程式）から未来を予想します。第6巻の章をのぞいてみましょう。\par\vspace{2.5mm}
\noindent\bknextcard{6-2}{表計算で解こう：オイラー法}{80 ℃のお湯は、10分後に何℃？}{短い時間なら、かけ算で近似できる（5-8）。それを表計算で何回もくり返す。}\hfill
\bknextcard{6-3}{積分するだけで解ける形}{落ちるボールの位置を、式で求める}{積分は微分の逆（5-11）と、積分定数は「はじめの状態」（5-13）。}\par\vspace{1.8mm}
\def\bknextrow{2}\noindent\bknextcard{6-4}{自分に比例してへる：冷めるお湯}{熱いほど速く冷める。答えは指数関数}{微分しても形が変わらない数 $e$（5-14）。答えが指数関数になる。}\hfill
\bknextcard{6-6}{ばねと振り子}{小さなゆれなら、周期はゆれの大きさによらない}{$\sin\cdot\cos$ の微分（5-15）と、$\sin\theta\approx\theta$ の近似（5-17）。}\par\vspace{1.8mm}
\def\bknextrow{3}\noindent\bknextcard{6-10}{オイラーの公式}{指数関数と三角関数が、1本の式でつながる}{$e^x$、$\sin x$、$\cos x$ のテイラー展開（5-17）をならべてくらべる。}\hfill
\bknextcard{6-13}{ゆらし続けると：強制振動と共振}{小さな力でも、タイミングが合うと大きくゆれる}{いちばん大きくなる所を、微分が0の点で見つける（5-7）。}\par
\bkfukidashi{delta}{第5巻の道具は、この先もずっと使うよ。チェック表で自信のない章を見直してから、第6巻へ進もう！}

\clearpage
\bkbkhead{この巻の道具は、この先でこう使う}
{\small 第5巻の章と、その道具を使う第6巻・第9巻の章です（章のはじめの「道しるべ」の「この章の先に」から集めました）。}\par\vspace{3mm}
\begin{bkcenter}
\begin{tblr}{colspec={Q[l,m,wd=52mm] Q[l,m,wd=80mm]}, row{1}={bg=bkVol!14, font=\sffamily\bfseries\footnotesize}, hlines={0.4pt,bkLine}, rowsep=1.5pt, cells={font=\footnotesize}}
この巻の章 & この道具を使う章 \\
{\bknextpill{5-5}\ 位置→速度→加速度} & {\bknextpill{6-1}\ 微分方程式ってなに？\quad \bknextpill{6-15}\ まとめ：物理は「微分方程式＋はじめの状態」で書ける} \\
{\bknextpill{5-7}\ いちばん高いのはどこ？：最大と最小} & {\bknextpill{6-13}\ ゆらし続けると：強制振動と共振} \\
{\bknextpill{5-8}\ 短い時間なら、かけ算で近似できる} & {\bknextpill{6-2}\ 表計算で解こう：オイラー法} \\
{\bknextpill{5-11}\ 積分は微分の逆} & {\bknextpill{6-3}\ 積分するだけで解ける形} \\
{\bknextpill{5-12}\ 物理での積分：仕事と流量} & {\bknextpill{6-15}\ まとめ：物理は「微分方程式＋はじめの状態」で書ける} \\
{\bknextpill{5-13}\ 積分定数は「はじめの状態」} & {\bknextpill{6-1}\ 微分方程式ってなに？\quad \bknextpill{6-3}\ 積分するだけで解ける形} \\
{\bknextpill{5-14}\ ふしぎな数 e} & {\bknextpill{6-4}\ 自分に比例してへる：冷めるお湯\quad \bknextpill{9-5}\ 微分をかけ算に変える翻訳機：ラプラス変換} \\
{\bknextpill{5-15}\ sin・cosを微分すると？} & {\bknextpill{6-6}\ ばねと振り子\quad \bknextpill{6-10}\ オイラーの公式} \\
{\bknextpill{5-16}\ 積分の道具箱：置換積分と部分積分} & {\bknextpill{6-4}\ 自分に比例してへる：冷めるお湯\quad \bknextpill{9-1}\ 波をたし合わせて形を作る：フーリエ級数\quad \bknextpill{9-5}\ 微分をかけ算に変える翻訳機：ラプラス変換} \\
{\bknextpill{5-17}\ 近くでは多項式で近似できる：テイラー展開} & {\bknextpill{6-6}\ ばねと振り子\quad \bknextpill{6-10}\ オイラーの公式\quad \bknextpill{9-1}\ 波をたし合わせて形を作る：フーリエ級数} \\
{\bknextpill{5-18}\ まとめ：微分したときのふるまい} & {\bknextpill{6-1}\ 微分方程式ってなに？\quad \bknextpill{6-4}\ 自分に比例してへる：冷めるお湯\quad \bknextpill{6-6}\ ばねと振り子} \\
\end{tblr}
\end{bkcenter}

%% ---------------- さくいん ----------------
\bkbacktitle{さくいん}
{\small 数字はページです。その語がはじめて出てくるページや、くわしく説明しているページをのせています。}\par\vspace{2mm}
\begin{multicols}{2}\small\setlength{\parskip}{0pt}
\directlua{bk.printidx()}
\end{multicols}

%% ---------------- 奥付 ----------------
\bkpadfour
\global\bk@decofalse\bkfolio{false}
\null\vfill
\noindent\begin{minipage}{\linewidth}\small
{\sffamily\bfseries\Large\color{bkDeep}物理を理解するための数学　第5巻}\par\vspace{1mm}
{\sffamily\bfseries\large 変化をつかむ　微分と積分}\par\vspace{3mm}
{\color{bkVol}\hrule height 1pt}\vspace{3mm}
\begin{tabular}{@{}l@{\hspace{6mm}}l@{}}
編者 & 「物理を理解するための数学」編集委員会 \\
 & 委員長　小学校長（まとめ役） \\
 & 委員　数学・電気工学・機械工学・教育学の大学教授、幼稚園教諭、\\
 & \phantom{委員　}特別支援学校教諭、心理士、科学解説者、小学校教諭 \\
デザイン & ユニバーサルデザイン担当委員 \\
挿絵 & イラスト担当委員 \\
版 & 見本版（第6回編集委員会で修正）　2026年10月 \\
\end{tabular}\par\vspace{3mm}
{\color{bkVol}\hrule height 1pt}\vspace{2mm}
{\footnotesize この本の問題の答えは、巻末の「こたえ」にあります。くわしい解き方は、別冊の問題集と指導書にのせます。}
\end{minipage}
\makeatother
